\documentclass[11pt]{article}

\usepackage[margin=0.85in]{geometry}
\usepackage{amsmath}
\usepackage{amssymb}
\usepackage{array}
\usepackage{booktabs}
\usepackage{hyperref}
\usepackage{longtable}
\usepackage{xcolor}
\usepackage{listings}
\usepackage{microtype}

\hypersetup{
  colorlinks=true,
  linkcolor=blue!50!black,
  urlcolor=blue!50!black,
  pdftitle={AthenaK Native Multigrid Tutorial},
  pdfauthor={Cursor agent}
}

\lstset{
  basicstyle=\ttfamily\small,
  columns=fullflexible,
  breaklines=true,
  frame=single,
  framerule=0.2pt,
  rulecolor=\color{black!30}
}

\newcommand{\file}[1]{\nolinkurl{#1}}
\newcommand{\code}[1]{\texttt{#1}}

\title{AthenaK Native Multigrid Tutorial}
\author{Concepts and code-reading guide}
\date{May 21, 2026}

\begin{document}
\maketitle

\section{What Problem Multigrid Solves}

Multigrid is an algorithm for solving elliptic equations, such as Poisson-like
equations or the CTS constraint equations.  These problems are difficult because
local relaxation methods remove short-wavelength errors quickly, but long-
wavelength errors decay slowly.

For a simple scalar problem,
\[
  L u = f,
\]
where \(L\) is a discretized differential operator, the solver tracks:
\begin{itemize}
  \item \(u\): the current solution estimate.
  \item \(f\): the right-hand side or source.
  \item \(r = f - L(u)\): the residual or defect.
  \item \(e = u_{\rm exact}-u\): the unknown error.
\end{itemize}

If \(u\) is wrong by a smooth, slowly varying error, pointwise smoothing on the
fine grid sees only a tiny local slope and converges slowly.  On a coarser grid,
the same smooth error is represented by fewer cells and becomes easier to remove.
The multigrid idea is therefore:
\begin{enumerate}
  \item Smooth on the fine grid to remove oscillatory errors.
  \item Restrict the remaining smooth residual/error information to a coarser
        grid.
  \item Correct the coarse-grid problem.
  \item Prolongate the coarse correction back to the fine grid.
  \item Smooth again to remove interpolation noise.
\end{enumerate}

\section{A Minimal V-Cycle}

The classic V-cycle looks like this:

\begin{lstlisting}
V(level):
  if level is coarsest:
    solve or smooth many times
  else:
    pre-smooth u_level
    residual_level = f_level - L_level(u_level)
    restrict residual and solution to level-1
    V(level-1)
    prolongate coarse correction to level
    post-smooth u_level
\end{lstlisting}

The ``V'' name comes from the path through the levels:

\begin{center}
\begin{tabular}{ccccccccc}
fine & $\rightarrow$ & medium & $\rightarrow$ & coarse & $\rightarrow$ & medium & $\rightarrow$ & fine \\
smooth & & restrict & & solve & & prolong & & smooth
\end{tabular}
\end{center}

\section{Linear Correction vs FAS}

For a linear problem \(Lu=f\), one often solves the error equation
\[
  L e = r.
\]
For nonlinear systems, such as the CTS residual \(F(u)=0\), AthenaK uses the
full approximation storage (FAS) pattern.  In FAS, coarse levels store a full
approximation \(u_H\), not just an error correction.  The coarse source is adjusted
so that the coarse problem is consistent with the fine-grid solution:
\[
  F_H(u_H) = f_H.
\]

In the current code, the physics-specific class supplies:
\begin{itemize}
  \item a residual/operator evaluation,
  \item a smoother,
  \item a defect computation,
  \item an FAS right-hand-side update.
\end{itemize}

The generic multigrid driver supplies hierarchy traversal, restriction,
prolongation, boundary exchange, root-grid handling, and convergence norms.

\section{The Main Code Objects}

\begin{longtable}{p{0.32\linewidth} p{0.62\linewidth}}
\toprule
\textbf{File/class} & \textbf{Purpose} \\
\midrule
\endfirsthead
\toprule
\textbf{File/class} & \textbf{Purpose} \\
\midrule
\endhead
\file{src/multigrid/multigrid.hpp} \newline \code{Multigrid} &
Owns per-level arrays on MeshBlocks or on the root grid.  Provides generic
operations such as loading data, restricting, prolongating, computing norms, and
storing old solutions. \\
\addlinespace
\file{src/multigrid/multigrid.hpp} \newline \code{MGOctet} &
Small \(2^3\)-interior cell object used to bridge static mesh-refinement
coarse/fine boundaries.  It stores \code{u}, \code{src}, \code{def},
\code{uold}, and optional coefficients. \\
\addlinespace
\file{src/multigrid/multigrid\_driver.cpp} \newline \code{MultigridDriver} &
Controls the global solve: builds the hierarchy, chooses current levels, runs
V-cycles/FMG cycles, moves data between MeshBlocks/root/octets, and computes
global norms with MPI reduction. \\
\addlinespace
\file{src/multigrid/multigrid\_tasks.cpp} &
Builds task lists for boundary exchange, smoothing, restriction, prolongation,
fine/coarse ghost filling, and send/receive cleanup. \\
\addlinespace
\file{src/z4c/id\_solve.hpp} \newline \code{IDCTSMultigrid} &
CTS-specific \code{Multigrid} subclass.  Implements \code{SmoothPack},
\code{CalculateDefectPack}, and \code{CalculateFASRHSPack}. \\
\addlinespace
\file{src/z4c/id\_solve.cpp} \newline \code{IDCTSMultigridDriver} &
CTS-specific driver.  Sets solver parameters, coefficient count, stencil choices,
loads CTS free data, calls generic multigrid cycles, and applies the solution. \\
\bottomrule
\end{longtable}

\section{The Arrays on Each Level}

Each \code{Multigrid} object allocates one array per multigrid level:

\begin{center}
\begin{tabular}{lll}
\toprule
\textbf{Array} & \textbf{Meaning} & \textbf{CTS example} \\
\midrule
\code{u\_} & current solution approximation & \(\delta\Psi,\delta\beta^i\) \\
\code{src\_} & source/RHS for the level & zero on finest CTS level; FAS RHS on coarser levels \\
\code{def\_} & defect/residual storage & \(src-F(u)\) \\
\code{coeff\_} & problem coefficients/free data & conformal metric, \(K\), \(\alpha\), sources \\
\code{uold\_} & saved solution before coarse solve & used for FAS correction \\
\bottomrule
\end{tabular}
\end{center}

The constructor \code{Multigrid::Multigrid()} creates either:
\begin{itemize}
  \item MeshBlock-local levels when given a \code{MeshBlockPack}, or
  \item root-grid levels when given \code{nullptr}.
\end{itemize}

MeshBlock levels require cubic, power-of-two MeshBlocks.  Root levels come from
the number of root MeshBlocks in each direction.  Each coarser level halves the
number of active cells per direction.

\section{Levels in AthenaK}

There are three kinds of levels in this implementation:

\subsection{MeshBlock Levels}

These are the normal levels inside each MeshBlock.  For a \(64^3\) MeshBlock,
the internal multigrid levels are \(64^3,32^3,16^3,\ldots,1^3\), plus ghost
zones.  \code{mglevels\_} owns these arrays.

\subsection{Root Grid Levels}

The root grid is built from the logical arrangement of root MeshBlocks.  For a
\(2\times2\times2\) root MeshBlock layout, the root-grid active region starts
as \(2^3\), then can coarsen to \(1^3\).  \code{mgroot\_} owns this part.

\subsection{Octets for Mesh Refinement}

For static mesh refinement, an octet represents a parent cell with \(2\times2
\times2\) refined children.  Octets let the V-cycle pass through refinement
levels without requiring every part of the domain to be uniformly refined.

For the native CTS solver, there are explicit guards:
\begin{itemize}
  \item adaptive AMR is not supported,
  \item sixth-order CTS on SMR is not supported by the current octet bridge.
\end{itemize}

\section{Setup Before a Solve}

The generic setup path is in \code{MultigridDriver::PrepareForAMR()} and
\code{MultigridDriver::SetupMultigrid()}.

\subsection{\code{PrepareForAMR()}}

Despite the name, this runs for every multigrid solve.  It:
\begin{itemize}
  \item records the root mesh level,
  \item detects whether the MeshBlock layout changed,
  \item recomputes how many refinement levels exist,
  \item reallocates MeshBlock multigrid arrays if needed,
  \item rebuilds octets if SMR/AMR levels exist,
  \item updates per-block cell widths.
\end{itemize}

\subsection{\code{SetupMultigrid()}}

This computes:
\[
  n_{\rm total} = n_{\rm root} + n_{\rm MeshBlock} + n_{\rm refinement} - 1.
\]
It sets \code{current\_level\_} to the finest level and initializes the FMG level
bookkeeping.

For CTS, \code{IDCTSMultigridDriver::Solve()} wraps setup with the physics data:
\begin{enumerate}
  \item build CTS free data,
  \item load \code{u\_cts} into \code{u\_},
  \item load zero source,
  \item load free-data coefficients and fill their finest-level halos,
  \item restrict coefficients and fill coefficient halos on each coarse level,
  \item setup the hierarchy,
  \item transfer coefficients to root and octets.
\end{enumerate}

\section{Restriction}

Restriction moves information from a fine grid to a coarser grid.  In a V-cycle,
restriction is used on the way down.

\code{Multigrid::RestrictPack()} does:
\begin{enumerate}
  \item call the physics-specific \code{CalculateDefectPack()},
  \item restrict \code{def\_} to the coarser \code{src\_},
  \item restrict \code{u\_} to the coarser \code{u\_},
  \item decrement \code{current\_level\_}.
\end{enumerate}

For nonlinear FAS, restricting the current solution as well as the defect is
important because the coarse grid stores a full approximation, not just an error.

In the octet path, \code{RestrictOctets()} computes each octet defect and writes
the restricted defect/source and restricted solution into either a coarser octet
or the root grid.

\section{Prolongation and Correction}

Prolongation moves a coarse-grid correction back to a finer grid.  AthenaK uses:
\begin{itemize}
  \item \code{ProlongateAndCorrectPack()} for normal V-cycle corrections,
  \item \code{FMGProlongatePack()} for direct FMG initialization,
  \item \code{ProlongateAndCorrectOctets()} for octet coarse/fine paths.
\end{itemize}

The normal V-cycle path first computes a FAS correction:
\[
  \Delta u_H = u_H - u_{H,{\rm old}},
\]
then prolongates \(\Delta u_H\) and adds it to the next-finer solution.

For MeshBlock-local levels, the code comments describe this as trilinear
correction.  For some root/octet bridge paths, the implementation uses a
3-by-3-by-3 coarse neighborhood with cubic-like weights.

\section{Smoothing}

Smoothing is the problem-specific local relaxation.  The generic driver only
knows it can call:
\[
  \code{pmg->SmoothPack(color)}.
\]

For CTS, \code{IDCTSMultigrid::SmoothPack()} dispatches to
\code{SmoothImpl<2/3/4>()} based on the effective CTS stencil.  The finest
MeshBlock MG level uses the Z4c stencil, while coarser MeshBlock/root levels use
\code{id\_solve/mg\_coarse\_fd\_stencil} (default \code{2} for higher-order fine
runs).  The smoothing kernel is red-black colored:
\[
  c = (\mathrm{color}+i+j+k)\bmod 2.
\]
Cells of one color are updated first, boundary/ghost zones are refreshed, then
the other color is updated.  CTS now uses
\code{id\_solve/smoother\_update=frozen\_view}: before each smoothing visit the
active level is copied into a dedicated frozen scratch view.  Operator
evaluations, diagonal updates, Newton-GS Jacobians, and line-search trial
residuals read from the frozen view while accepted updates write to the live
solution.  This removes same-color read/write races for mixed and high-order
stencils without expanding the task list to an 8/27/64-color schedule.

The default CTS update is:
\[
  u_v \leftarrow u_v - \omega\,
  \frac{F_v(u)-src_v}{D_v},
\]
where \(D_v\) is an approximate diagonal of the CTS operator.  This is an
under-relaxed diagonal nonlinear smoother.

The CTS driver can also use \code{id\_solve/smoother = newton\_gs}.  In that
mode, each active red/black point forms the four-component local residual for
\((\Psi,\beta^x,\beta^y,\beta^z)\), finite-differences a scale-aware local
\(4\times4\) Jacobian of the CTS operator, solves the coupled Newton correction,
and applies a residual-checked damped update.  If the local linear solve is
singular, ill-conditioned, non-finite, or cannot produce a local residual
decrease after backtracking, the code falls back to the diagonal update for that
point.
This coupled update is the better long-term algorithmic target for nonlinear
FAS, but it is not automatically faster: it does more work at every active
point.  Production comparisons should therefore use defect reduction per
wall-clock time, not stable \(\omega\) alone.

Why red-black?  A centered stencil couples nearest neighbors.  If the grid is
colored like a checkerboard, red cells depend mostly on black neighbors and vice
versa.  Updating one color at a time gives Gauss-Seidel-like in-place relaxation
while preserving parallelism within each color.  For CTS high-order and mixed
derivative operators, the frozen view is what makes this two-color schedule
race-free.

\section{The Current V-Cycle in Code}

The V-cycle entry point is:
\[
  \code{MultigridDriver::SolveVCycle(pdriver, npresmooth, npostsmooth)}.
\]

It does:
\begin{lstlisting}
startlevel = current_level
coffset_ ^= 1
while current_level > 0:
  OneStepToCoarser(npresmooth)
SolveCoarsestGrid()
while current_level < startlevel:
  OneStepToFiner(npostsmooth)
\end{lstlisting}

\subsection{One Step to Coarser}

\code{OneStepToCoarser()} depends on which kind of level is active:
\begin{itemize}
  \item MeshBlock level: run the \code{mg\_to\_coarser} task list.
  \item Octet level: smooth octets, refresh octet boundaries, restrict octets.
  \item Root level: apply root boundary, optionally compute FAS RHS, smooth
        red/black with root boundary after each color, restrict root grid.
\end{itemize}

The MeshBlock task list includes:
\begin{enumerate}
  \item start receive,
  \item send boundary data,
  \item receive/unpack boundary data,
  \item physical boundary fill,
  \item fine/coarse ghost fill,
  \item FAS RHS computation,
  \item red smoothing,
  \item boundary exchange/fill,
  \item black smoothing,
  \item boundary exchange/fill,
  \item optional extra red/black pair if \code{nsmooth > 1},
  \item restriction,
  \item clear send/recv.
\end{enumerate}

\subsection{Coarsest Grid Solve}

\code{SolveCoarsestGrid()} does not use a direct matrix solve.  It applies root
boundaries, stores old data, computes FAS RHS, then performs a number of
red-black smoothing passes.  The number of passes is roughly the size of the
coarsest root grid in one dimension.  If the coarsest root grid is \(1^3\) and
average subtraction is enabled, it may zero-clear instead.

\subsection{One Step to Finer}

The upward path calls \code{OneStepToFiner()}.  It prolongates a coarse
correction and then performs post-smoothing.  On MeshBlock levels, the
\code{mg\_to\_finer} task list performs any needed boundary communication before
prolongation, then boundary communication and red-black post-smoothing after
prolongation.  On the final step to the finest level, it performs a final
boundary exchange so the defect norm sees up-to-date ghost zones.

\section{Full Multigrid (FMG)}

FMG is an optional startup strategy controlled by \code{full\_multigrid}.  Instead
of beginning on the finest grid with a rough guess, FMG:
\begin{enumerate}
  \item restricts the problem down to very coarse levels,
  \item solves/coarsely smooths there,
  \item prolongates to a finer level,
  \item performs one or more V-cycles,
  \item repeats until the finest level is reached,
  \item then continues with the normal MG solve.
\end{enumerate}

In code, this is \code{SolveFMG()}, \code{SolveFMGCoarser()}, and
\code{FMGProlongate()}.  CTS has \code{id\_solve/full\_multigrid} and
\code{id\_solve/fmg\_ncycle} controls.

\section{Boundary Handling}

Boundary handling is essential because each smoother and residual evaluation
uses ghost cells.

\subsection{MeshBlock Boundary Exchange}

MeshBlock task lists use:
\begin{itemize}
  \item \code{StartReceive()} \(\to\) \code{InitRecvMG()},
  \item \code{SendBoundary()} \(\to\) \code{PackAndSendMG()},
  \item \code{RecvBoundary()} \(\to\) \code{RecvAndUnpackMG()},
  \item \code{ClearSend()} and \code{ClearRecv()} for cleanup.
\end{itemize}

These are scheduled around smoothing, restriction, and prolongation so stencil
operations see valid ghost data.

CTS also exchanges and fills coefficient/free-data halos.  After
\code{LoadCoefficients()} and after each \code{RestrictCoefficients()} step, the
coefficient arrays use the same same-level MG exchange buffers as \code{u\_}.
Their physical boundaries are filled as free data: periodic faces wrap and
non-periodic faces use zero-gradient copies.  When refinement is present, the
coefficient arrays then call \code{FillFineCoarseMGGhosts()} so CTS residuals
near fine/coarse interfaces do not read stale free-data ghosts.  They do not use
correction-variable odd reflection from \code{mg\_zerofixed}.

\subsection{Physical Boundaries}

The base multigrid driver initializes multigrid physical BCs from mesh BCs:
\begin{itemize}
  \item periodic mesh faces stay periodic,
  \item non-periodic mesh faces become \code{mg\_zerofixed} in the generic base
        driver.
\end{itemize}

\code{mg\_zerofixed} uses odd reflection of the correction variable.  This
approximately enforces zero value at the physical boundary.  \code{mg\_zerograd}
uses even reflection.  Multipole BC support exists in the base MG machinery and
is used by gravity.  The CTS driver overrides the non-periodic default with
\code{id\_solve/mg\_bc}; its current default is \code{zerograd}, with
\code{zerofixed} available for older behavior.

\subsection{Fine/Coarse Ghosts}

When SMR/AMR is present, \code{FillFCBoundary()} calls
\code{FillFineCoarseMGGhosts()}.  For root/octet paths, octet boundary helpers
copy same-level data or interpolate from coarser neighbors.  CTS also transfers
coefficient/free-data arrays through this hierarchy.

The current octet coefficient bridge still uses the generic compact
coarse-buffer path, not the full Z4c high-order prolongation kernels.  That is
why sixth-order CTS on SMR remains guarded off even though unigrid high-order CTS
operators are supported.  Octet CTS operators therefore default to
\code{id\_solve/mg\_coarse\_fd\_stencil}; \code{id\_solve/octet\_fd\_stencil}
is an explicit override.

\section{Residual and Norms}

\code{CalculateDefectPack()} is physics-specific.  For CTS it evaluates
\[
  def_v = src_v - F_v(u).
\]
\code{Multigrid::CalculateDefectNorm()} then reduces this defect over active
cells.  For L1 and L2 norms the code multiplies by cell volume.  The driver
\code{MultigridDriver::CalculateDefectNorm()} performs an MPI all-reduce and
divides by the total domain volume; L2 norms are square-rooted:
\[
  \|def\|_2 =
  \left(\frac{1}{V}\int |def|^2\,dV\right)^{1/2}.
\]

For multi-variable systems like CTS, the CTS driver computes each variable norm
and then reports the total as the Euclidean combination:
\[
  \|def\|_{\rm total}
  = \sqrt{\sum_v \|def_v\|_2^2}.
\]

\section{How CTS Uses the Generic Multigrid}

The CTS solve is a good concrete example:
\begin{enumerate}
  \item \code{IDConformalThinSandwich::BuildFreeData()} builds all coefficients.
  \item \code{IDCTSMultigridDriver::Solve()} loads the initial guess, source, and
        coefficients into the multigrid arrays.
  \item Generic V-cycle machinery calls back into:
        \code{SmoothPack()}, \code{CalculateDefectPack()}, and
        \code{CalculateFASRHSPack()}.
  \item \code{SmoothPack()} calls the selected CTS smoother
        (\code{diagonal} by default, or \code{newton\_gs}).
  \item \code{CalculateDefectPack()} evaluates the CTS residual.
  \item \code{CalculateFASRHSPack()} adds the CTS operator value to the coarse
        FAS source.
  \item After convergence or fixed V-cycles, \code{RetrieveResult()} copies the
        finest \code{u} back to \code{u\_cts}.
  \item \code{ApplySolution()} reconstructs ADM/Z4c data.
\end{enumerate}

\section{Fixed-Iteration vs Threshold Mode}

The generic driver and CTS driver support two normal modes:

\begin{itemize}
  \item \textbf{Threshold mode}: if \code{threshold >= 0}, V-cycles continue until
        the defect drops below the threshold or the iteration cap is reached.
  \item \textbf{Fixed-count mode}: if \code{threshold < 0}, the solver runs
        \code{niteration} V-cycles.
\end{itemize}

CTS adds best-solution guards:
\begin{itemize}
  \item \code{keep\_best\_solution}: stores the best-so-far correction if later
        V-cycles diverge.
  \item \code{stop\_on\_defect\_increase}: optionally stops when the defect grows
        beyond a tolerance.
  \item \code{reject\_worse}: rejects the solved update if the final defect is
        worse than the initial defect.
\end{itemize}

\section{How to Read the Multigrid Code}

A useful reading order is:
\begin{enumerate}
  \item \file{src/multigrid/multigrid.hpp}: learn the arrays, levels, \code{MGOctet},
        and virtual physics hooks.
  \item \file{src/multigrid/multigrid.cpp}: read \code{LoadFinestData},
        \code{LoadSource}, \code{LoadCoefficients}, \code{RestrictPack},
        \code{ProlongateAndCorrectPack}, and \code{CalculateDefectNorm}.
  \item \file{src/multigrid/multigrid\_driver.cpp}: read \code{SetupMultigrid},
        \code{SolveVCycle}, \code{SolveCoarsestGrid}, \code{SolveFMG}, and
        \code{MGRootBoundary}.
  \item \file{src/multigrid/multigrid\_tasks.cpp}: read task-list construction
        for \code{mg\_to\_coarser} and \code{mg\_to\_finer}.
  \item \file{src/z4c/id\_solve.cpp}: read how CTS plugs into the generic hooks.
\end{enumerate}

\section{Mental Model for Debugging}

When multigrid misbehaves, classify the failure:
\begin{itemize}
  \item \textbf{Smoother instability}: defect grows during V-cycles.  Try smaller
        \(\omega\), compare \code{smoother=diagonal} with
        \code{smoother=newton\_gs}, inspect the local Jacobian/diagonal
        approximation, or use the best-solution guard.
  \item \textbf{Boundary/communication bug}: MPI stalls or rank-dependent results.
        Inspect task dependencies around \code{StartReceive}, \code{SendBoundary},
        \code{RecvBoundary}, \code{ClearSend}, and \code{ClearRecv}.
  \item \textbf{Coarse-grid inconsistency}: fine-grid residual improves briefly
        but correction from coarse levels hurts.  Inspect restriction,
        prolongation, and coefficient restriction.
  \item \textbf{Stencil/halo mismatch}: high-order operators use invalid ghost
        data or silently reduce order.  Check stencil radius, ghost count, and
        whether any precomputed free-data fields need their own ghost halo.
  \item \textbf{Bad free data or source scaling}: initial defect changes strongly
        with resolution or physical constraints do not improve after solving.
\end{itemize}

For CTS specifically, the diagonal smoother remains the conservative default and
is simple and parallel.  The optional \code{newton\_gs} smoother includes the
local coupling among \(\Psi\) and \(\beta^i\), but still needs damping and is
more expensive because it finite-differences and solves a \(4\times4\) system at
each active point.  Current production safeguards include a residual-checked
line search, scale-aware local linear algebra, finite-value guards, update
limits, and MPI-reduced smoother counters.  Compare smoothers by defect
reduction per wall-clock time rather than by stable \(\omega\) alone.  In the
current \(64^3\), second-order, 40-cycle critical-Kerr comparison,
\code{newton\_gs} reached a lower final defect than \code{diagonal} at
\(\omega=0.05\), but was still slower per wall-clock time.

\section{Glossary}

\begin{longtable}{p{0.25\linewidth} p{0.67\linewidth}}
\toprule
\textbf{Term} & \textbf{Meaning} \\
\midrule
\endfirsthead
\toprule
\textbf{Term} & \textbf{Meaning} \\
\midrule
\endhead
Smoother & Local relaxation step that damps high-frequency error. \\
Residual/defect & Difference between desired equation and current operator value,
\(f-L(u)\) or \(src-F(u)\). \\
Restriction & Fine-to-coarse transfer. \\
Prolongation & Coarse-to-fine transfer. \\
V-cycle & Down to coarse levels, coarse solve/smooth, back to fine levels. \\
FAS & Full Approximation Storage, nonlinear multigrid method where coarse grids
store full solutions rather than only corrections. \\
Root grid & Coarse grid built from the logical root MeshBlock layout. \\
MeshBlock level & Multigrid hierarchy inside each MeshBlock. \\
Octet & \(2^3\)-interior-cell structure that bridges refinement levels. \\
Ghost cells & Extra boundary cells needed for stencils and inter-block exchange. \\
\bottomrule
\end{longtable}

\end{document}
